\documentclass[11pt, a4paper]{article}
\usepackage[utf8]{inputenc}
\usepackage[margin=0.75in]{geometry}
\usepackage{amsmath, amssymb}
\usepackage{enumitem}
\usepackage{multicol}

\title{\textbf{NEET Biology: ETS and Oxidative Phosphorylation — High-Yield Difficult Questions}}
\author{Target: NEET UG}
\date{}

\begin{document}

\maketitle

\section*{Questions}

\begin{enumerate}[label=\textbf{\arabic*.}]

\item Which of the following correctly arranges the four multi-protein complexes of the mitochondrial Electron Transport System (ETS) along with their exact biochemical designations according to NCERT?
\begin{enumerate}[label=(\Alph*)]
    \item Complex I: Succinate dehydrogenase; Complex II: $\text{NADH}$ dehydrogenase; Complex III: Cytochrome $c$ oxidase; Complex IV: Cytochrome $bc_1$
    \item Complex I: $\text{NADH}$ dehydrogenase; Complex II: Succinate dehydrogenase; Complex III: Cytochrome $bc_1$ complex; Complex IV: Cytochrome $c$ oxidase
    \item Complex I: $\text{NADH}$ dehydrogenase; Complex II: Cytochrome $bc_1$ complex; Complex III: Succinate dehydrogenase; Complex IV: $\text{ATP}$ synthase
    \item Complex I: Cytochrome $c$ reductase; Complex II: $\text{FADH}_2$ oxidase; Complex III: Cytochrome $a-a_3$; Complex IV: Succinate-CoQ reductase
\end{enumerate}

\item Identify the two **mobile electron carriers** in the mitochondrial ETS that shuttle electrons between the stationary transmembrane complexes, along with their exact physical locations:
\begin{enumerate}[label=(\Alph*)]
    \item Ubiquinone (hydrophobic, within the inner mitochondrial membrane) and Cytochrome $c$ (small peripheral protein attached to the outer surface of the inner mitochondrial membrane)
    \item Ubiquinone (hydrophilic, in the intermembrane space) and Cytochrome $b$ (embedded in the outer mitochondrial membrane)
    \item Cytochrome $c$ (in the mitochondrial matrix) and Ferredoxin (within the inner mitochondrial membrane)
    \item Plastoquinone (in the inner membrane) and Cytochrome $a_3$ (on the matrix side of the inner membrane)
\end{enumerate}

\item \textbf{Assertion (A):} Oxidation of one molecule of $\text{NADH} + \text{H}^+$ yields $3\text{ molecules}$ of $\text{ATP}$, whereas oxidation of one molecule of $\text{FADH}_2$ yields only $2\text{ molecules}$ of $\text{ATP}$ in the mitochondrial ETS (per NCERT stoichiometry). \\
\textbf{Reason (R):} Electrons from $\text{FADH}_2$ enter the ETS via Complex II (Succinate dehydrogenase) directly to ubiquinone, thereby bypassing the proton-pumping Complex I ($\text{NADH}$ dehydrogenase).
\begin{enumerate}[label=(\Alph*)]
    \item Both (A) and (R) are true and (R) is the correct explanation of (A).
    \item Both (A) and (R) are true but (R) is not the correct explanation of (A).
    \item (A) is true but (R) is false.
    \item Both (A) and (R) are false.
\end{enumerate}

\item Which of the following prosthetic groups and metal centers are integral components of **Complex IV (Cytochrome $c$ oxidase)**?
\begin{enumerate}[label=(\Alph*)]
    \item $\text{FMN}$ and Iron-Sulfur ($\text{Fe}-\text{S}$) clusters
    \item $\text{FAD}$ and Heme $b_{560}$
    \item Cytochromes $a$ and $a_3$, and two copper centers ($\text{Cu}_A$ and $\text{Cu}_B$)
    \item Cytochromes $b_L$, $b_H$, and Rieske $\text{Fe}-\text{S}$ protein
\end{enumerate}

\item Why is the energy-coupling process in mitochondria specifically termed **oxidative phosphorylation**, whereas the analogous process in chloroplasts is termed **photophosphorylation**?
\begin{enumerate}[label=(\Alph*)]
    \item Because mitochondria synthesize $\text{ATP}$ using molecular oxygen directly as a phosphate donor, while chloroplasts use photons.
    \item Because in mitochondria, the energy of oxidation-reduction reactions creates the proton gradient required for $\text{ATP}$ synthesis, whereas in chloroplasts, light energy creates the proton gradient.
    \item Because oxidative phosphorylation occurs exclusively in the cytosol, while photophosphorylation occurs in the stroma.
    \item Because oxidative phosphorylation consumes $\text{ATP}$ to oxidize glucose, while photophosphorylation produces $\text{ATP}$ without a proton gradient.
\end{enumerate}

\item Match Column-I (ETS Inhibitor / Uncoupler) with Column-II (Specific Target / Mechanism of Action):
\begin{center}
\begin{tabular}{|l|l|}
\hline
\textbf{Column-I} & \textbf{Column-II} \\
\hline
P. Cyanide ($\text{CN}^-$) and Carbon monoxide ($\text{CO}$) & 1. Blocks proton channel ($\text{F}_0$ subunit) of $\text{ATP}$ synthase \\
Q. Rotenone and Amytal & 2. Binds to $\text{Fe}^{3+}/\text{Cu}_B$ center of Complex IV, blocking electron transfer to $\text{O}_2$ \\
R. Oligomycin & 3. Inhibits electron transfer from $\text{Fe}-\text{S}$ centers of Complex I to Ubiquinone \\
S. 2,4-Dinitrophenol ($\text{2,4-DNP}$) & 4. Protonophore that dissipates the $\text{H}^+$ gradient without $\text{ATP}$ synthesis \\
\hline
\end{tabular}
\end{center}
\begin{enumerate}[label=(\Alph*)]
    \item P-2, Q-3, R-1, S-4
    \item P-3, Q-2, R-1, S-4
    \item P-2, Q-1, R-3, S-4
    \item P-4, Q-3, R-2, S-1
\end{enumerate}

\item According to Peter Mitchell's **Chemiosmotic Hypothesis**, what is the exact direction of proton ($\text{H}^+$) pumping during electron transport, and where do protons accumulate to establish the electrochemical gradient (proton-motive force)?
\begin{enumerate}[label=(\Alph*)]
    \item Pumped from the intermembrane space into the mitochondrial matrix; protons accumulate in the matrix.
    \item Pumped from the mitochondrial matrix across the inner membrane into the intermembrane space; protons accumulate in the intermembrane space (lowering its pH).
    \item Pumped from the cytosol across the outer membrane into the intermembrane space; protons accumulate in the cytosol.
    \item Pumped from the mitochondrial matrix into the cytosol via porins; protons accumulate in the cytosol.
\end{enumerate}

\item According to NCERT Class 11 Biology, for each $\text{ATP}$ molecule produced by Complex V ($\text{ATP}$ synthase), how many protons ($\text{H}^+$) pass through the $\text{F}_0$ channel from the intermembrane space back into the matrix down the electrochemical proton gradient?
\begin{enumerate}[label=(\Alph*)]
    \item $1\text{ H}^+$
    \item $2\text{ H}^+$
    \item $3\text{ H}^+$
    \item $4\text{ H}^+$
\end{enumerate}

\item Consider the structural and functional organization of **Complex V ($\text{ATP}$ synthase)**:
\begin{enumerate}[label=(\roman*)]
    \item The $\text{F}_1$ headpiece is a peripheral membrane protein complex protruding into the mitochondrial matrix and contains the catalytic sites for $\text{ATP}$ synthesis from $\text{ADP}$ and $\text{P}_i$.
    \item The $\text{F}_0$ component is an integral membrane protein complex embedded in the inner mitochondrial membrane that forms a transmembrane proton channel.
    \item Passage of protons through $\text{F}_0$ induces rotational conformational changes in the $\gamma$-subunit, driving the binding-change mechanism in the $\beta$-subunits of $\text{F}_1$.
    \item The $\text{F}_1$ subunit faces the intermembrane space, while $\text{F}_0$ is located on the outer mitochondrial membrane.
\end{enumerate}
Which of the above statements are correct?
\begin{enumerate}[label=(\Alph*)]
    \item (i), (ii), and (iii)
    \item (i) and (iv) only
    \item (ii), (iii), and (iv)
    \item (i), (ii), (iii), and (iv)
\end{enumerate}

\item What is the exact role of **molecular oxygen ($\text{O}_2$)** in aerobic cellular respiration?
\begin{enumerate}[label=(\Alph*)]
    \item It directly oxidizes glucose in the cytoplasm to release carbon dioxide.
    \item It acts as the final hydrogen/electron acceptor at Complex IV, removing electrons and protons from the system to form metabolic water ($\text{H}_2\text{O}$).
    \item It donates electrons to Complex I to initiate the electron transport chain.
    \item It binds to the $\text{F}_1$ subunit of $\text{ATP}$ synthase to catalyze the phosphorylation of $\text{ADP}$.
\end{enumerate}

\item How many molecules of **metabolic water ($\text{H}_2\text{O}$)** are formed at Complex IV when **one molecule of glucose** undergoes complete aerobic oxidation (assuming the malate-aspartate shuttle operates)?
\begin{enumerate}[label=(\Alph*)]
    \item $2\text{ H}_2\text{O}$
    \item $6\text{ H}_2\text{O}$
    \item $10\text{ H}_2\text{O}$
    \item $12\text{ H}_2\text{O}$
\end{enumerate}

\item Why does **Ubiquinone ($\text{CoQ}$)** play a unique central role in the mitochondrial electron transport chain compared to cytochromes?
\begin{enumerate}[label=(\Alph*)]
    \item Ubiquinone is a water-soluble protein located exclusively in the matrix.
    \item Ubiquinone can accept both electrons and protons (forming ubiquinol, $\text{QH}_2$) from Complexes I and II, and releases protons into the intermembrane space while transferring single electrons to Complex III (a proton-electron coupler).
    \item Ubiquinone transfers electrons directly to molecular oxygen without involving Complex IV.
    \item Ubiquinone contains a copper ion that undergoes reversible $\text{Cu}^{2+}/\text{Cu}^+$ reduction.
\end{enumerate}

\item If a student adds **2,4-Dinitrophenol ($\text{2,4-DNP}$)**—a classic chemical uncoupler—to actively respiring mitochondria supplied with succinate, $\text{ADP}$, $\text{P}_i$, and $\text{O}_2$, what will be the immediate physiological consequence?
\begin{enumerate}[label=(\Alph*)]
    \item Both oxygen consumption and $\text{ATP}$ synthesis will stop completely.
    \item Oxygen consumption and electron transport will continue (or even accelerate), while $\text{ATP}$ synthesis will drop to zero and free energy will be dissipated as heat.
    \item $\text{ATP}$ synthesis will double while oxygen consumption stops.
    \item Succinate will be converted back into glucose via reverse glycolysis.
\end{enumerate}

\item \textbf{Statement I:} In brown adipose tissue (BAT) of hibernating mammals and human newborns, non-shivering thermogenesis is mediated by **Thermogenin (UCP1)**, a natural uncoupling protein in the inner mitochondrial membrane. \\
\textbf{Statement II:} Thermogenin provides an alternative pathway for protons to re-enter the mitochondrial matrix without passing through $\text{ATP}$ synthase, releasing the energy of the proton gradient directly as heat.
\begin{enumerate}[label=(\Alph*)]
    \item Both Statement I and Statement II are correct.
    \item Both Statement I and Statement II are incorrect.
    \item Statement I is correct, but Statement II is incorrect.
    \item Statement I is incorrect, but Statement II is correct.
\end{enumerate}

\item In the ETS, electrons move spontaneously from carriers with more negative reduction potentials ($E^{\circ\prime}$) to carriers with more positive reduction potentials. Which of the following sequences represents the correct thermodynamic order of electron flow?
\begin{enumerate}[label=(\Alph*)]
    \item $\text{NADH} \rightarrow \text{FMN} \rightarrow \text{Fe-S} \rightarrow \text{Ubiquinone} \rightarrow \text{Cyt } b \rightarrow \text{Cyt } c_1 \rightarrow \text{Cyt } c \rightarrow \text{Cyt } a \rightarrow \text{Cyt } a_3 \rightarrow \text{O}_2$
    \item $\text{NADH} \rightarrow \text{FAD} \rightarrow \text{Cyt } c \rightarrow \text{Ubiquinone} \rightarrow \text{Cyt } b \rightarrow \text{Cyt } a_3 \rightarrow \text{Cyt } a \rightarrow \text{O}_2$
    \item $\text{FADH}_2 \rightarrow \text{FMN} \rightarrow \text{Cyt } a \rightarrow \text{Cyt } c \rightarrow \text{Ubiquinone} \rightarrow \text{Cyt } b \rightarrow \text{O}_2$
    \item $\text{NADH} \rightarrow \text{Ubiquinone} \rightarrow \text{FMN} \rightarrow \text{Cyt } c \rightarrow \text{Cyt } b \rightarrow \text{Cyt } a \rightarrow \text{O}_2$
\end{enumerate}

\item Which of the following complexes in the mitochondrial inner membrane do **not** pump protons ($\text{H}^+$) from the matrix to the intermembrane space?
\begin{enumerate}[label=(\Alph*)]
    \item Complex I only
    \item Complex II only
    \item Complexes I, III, and IV
    \item Complex III and Complex IV
\end{enumerate}

\item Because the inner mitochondrial membrane is impermeable to $\text{NADH}$, eukaryotic cells use shuttle systems to transfer reducing equivalents from cytosolic $\text{NADH}$ (generated during glycolysis) into the mitochondrion. Compare the **Malate-Aspartate shuttle** (heart, liver, kidney) with the **Glycerol-3-Phosphate shuttle** (skeletal muscle, brain):
\begin{enumerate}[label=(\Alph*)]
    \item The malate-aspartate shuttle transfers electrons to mitochondrial $\text{NAD}^+$ (yielding $3\text{ ATP}$ per cytosolic $\text{NADH}$), whereas the glycerol-3-phosphate shuttle transfers electrons to mitochondrial $\text{FAD}$ (yielding $2\text{ ATP}$ per cytosolic $\text{NADH}$).
    \item Both shuttles transfer electrons directly to Complex IV, yielding $1\text{ ATP}$ per cytosolic $\text{NADH}$.
    \item The glycerol-3-phosphate shuttle yields $38\text{ ATP}$ per glucose, whereas the malate-aspartate shuttle yields $36\text{ ATP}$ per glucose.
    \item Both shuttles transport intact $\text{NADH}$ molecules through porin channels in the inner membrane.
\end{enumerate}

\item How many total $\text{ATP}$ molecules are produced from **1 molecule of glucose** in human skeletal muscle cells operating via the **glycerol-3-phosphate shuttle** (using classical NCERT conversion factors: $1\text{ NADH} = 3\text{ ATP}$, $1\text{ FADH}_2 = 2\text{ ATP}$)?
\begin{enumerate}[label=(\Alph*)]
    \item $32\text{ ATP}$
    \item $36\text{ ATP}$
    \item $38\text{ ATP}$
    \item $40\text{ ATP}$
\end{enumerate}

\item In **Complex III (Cytochrome $bc_1$ complex)**, how does the **Q-cycle** overcome the mechanistic mismatch between ubiquinol ($\text{QH}_2$, a two-electron carrier) and cytochrome $c$ (a one-electron carrier)?
\begin{enumerate}[label=(\Alph*)]
    \item By converting cytochrome $c$ into a two-electron carrier using $\text{FMN}$.
    \item By bifurcating the two electrons of $\text{QH}_2$: one electron travels via the Rieske $\text{Fe}-\text{S}$ protein and cytochrome $c_1$ to cytochrome $c$, while the second electron recycles via cytochromes $b_L$ and $b_H$ to partially reduce another ubiquinone molecule.
    \item By releasing one electron into the intermembrane space as a free radical.
    \item By passing both electrons simultaneously to Complex IV without involving cytochrome $c$.
\end{enumerate}

\item Why does **cyanide ($\text{CN}^-$)** cause rapid death within minutes when ingested or inhaled?
\begin{enumerate}[label=(\Alph*)]
    \item It competitively inhibits hexokinase in erythrocytes.
    \item It binds irreversibly to the heme $a_3\text{--Cu}_B$ binuclear center of Cytochrome $c$ oxidase (Complex IV), halting electron transfer to $\text{O}_2$, collapsing the proton gradient, and stopping oxidative phosphorylation in all aerobic tissues.
    \item It dissolves the outer mitochondrial membrane.
    \item It converts all cellular $\text{ATP}$ into cyclic $\text{AMP}$.
\end{enumerate}

\item If **Oligomycin** (an inhibitor of the $\text{F}_0$ proton channel of $\text{ATP}$ synthase) is added to intact, coupled mitochondria, why does **electron transport and oxygen consumption** also stop shortly afterward?
\begin{enumerate}[label=(\Alph*)]
    \item Oligomycin directly binds to Cytochrome $c$ oxidase as well.
    \item Without proton re-entry through $\text{F}_0$, the proton gradient across the inner membrane builds up to a maximum thermodynamic back-pressure, preventing Complexes I, III, and IV from pumping any more protons.
    \item Oligomycin destroys all $\text{NAD}^+$ in the matrix.
    \item Oxygen cannot diffuse across the inner membrane when $\text{F}_0$ is closed.
\end{enumerate}

\item Suppose an experimenter adds **Oligomycin** to coupled mitochondria (halting both $\text{ATP}$ synthesis and $\text{O}_2$ consumption), and then subsequently adds **2,4-Dinitrophenol ($\text{2,4-DNP}$)** to the same preparation. What will happen?
\begin{enumerate}[label=(\Alph*)]
    \item Both $\text{O}_2$ consumption and $\text{ATP}$ synthesis will resume normally.
    \item $\text{O}_2$ consumption (electron transport) will resume rapidly because $\text{2,4-DNP}$ dissipates the proton back-pressure, but $\text{ATP}$ synthesis will remain completely blocked.
    \item Neither $\text{O}_2$ consumption nor $\text{ATP}$ synthesis will change; both remain zero.
    \item $\text{ATP}$ synthesis will resume, while $\text{O}_2$ consumption remains zero.
\end{enumerate}

\item Which of the following cytochromes in the mitochondrial ETS contains **iron** in its heme prosthetic group, whereas **Cytochrome $c$ oxidase** additionally requires **copper**?
\begin{enumerate}[label=(\Alph*)]
    \item All cytochromes ($b$, $c_1$, $c$, $a$, $a_3$) contain iron in a heme porphyrin ring, and undergo reversible $\text{Fe}^{3+} \leftrightarrow \text{Fe}^{2+}$ redox transitions.
    \item Only cytochrome $c$ contains iron; cytochromes $b$ and $a$ contain magnesium.
    \item Cytochromes $b$ and $c_1$ contain copper, while cytochromes $a$ and $a_3$ contain zinc.
    \item Cytochrome $c$ contains manganese instead of iron.
\end{enumerate}

\item How is **Cytochrome $c$** structurally anchored to the inner mitochondrial membrane compared to Cytochromes $b$, $c_1$, $a$, and $a_3$?
\begin{enumerate}[label=(\Alph*)]
    \item Cytochrome $c$ is a transmembrane protein spanning both the outer and inner membranes.
    \item Cytochrome $c$ is a small, water-soluble peripheral protein loosely bound via electrostatic interactions to the outer face (intermembrane space side) of the inner mitochondrial membrane, allowing it to diffuse laterally between Complex III and Complex IV.
    \item Cytochrome $c$ is covalently bound to the $\text{F}_1$ headpiece on the matrix side.
    \item Cytochrome $c$ is dissolved freely in the cytosol.
\end{enumerate}

\item \textbf{Assertion (A):} In the respiratory electron transport chain, the proton-motive force ($\Delta p$) consists of two components: a chemical pH gradient ($\Delta \text{pH}$) and an electrical membrane potential ($\Delta \psi$). \\
\textbf{Reason (R):} Pumping positively charged $\text{H}^+$ ions out of the matrix leaves the mitochondrial matrix negatively charged and alkaline relative to the positively charged, acidic intermembrane space.
\begin{enumerate}[label=(\Alph*)]
    \item Both (A) and (R) are true and (R) is the correct explanation of (A).
    \item Both (A) and (R) are true but (R) is not the correct explanation of (A).
    \item (A) is true but (R) is false.
    \item Both (A) and (R) are false.
\end{enumerate}

\item Consider the **Respiratory Balance Sheet** in NCERT. Which of the following assumptions are required to arrive at the theoretical net gain of **$38\text{ ATP}$** per glucose molecule?
\begin{enumerate}[label=(\roman*)]
    \item There is a sequential, orderly pathway functioning, with one substrate forming the next (glycolysis $\rightarrow$ TCA cycle $\rightarrow$ ETS).
    \item The $\text{NADH}$ synthesized in glycolysis is transferred into the mitochondria and undergoes oxidative phosphorylation.
    \item None of the intermediates in the pathway are utilized to synthesize any other compound (no anabolic withdrawal).
    \item Only glucose is being respired; no other alternative substrates enter the pathway at any intermediary stage.
\end{enumerate}
\begin{enumerate}[label=(\Alph*)]
    \item (i) and (ii) only
    \item (i), (iii), and (iv) only
    \item (ii), (iii), and (iv) only
    \item (i), (ii), (iii), and (iv)
\end{enumerate}

\item During the reduction of molecular oxygen to water at Complex IV, how many electrons are required to completely reduce **one molecule of $\text{O}_2$** without releasing toxic reactive oxygen species (such as superoxide $\text{O}_2^{\bullet-}$ or peroxide $\text{O}_2^{2-}$)?
\begin{enumerate}[label=(\Alph*)]
    \item $1\text{ electron}$
    \item $2\text{ electrons}$
    \item $4\text{ electrons}$ (along with $4\text{ H}^+$ from the matrix to form $2\text{ H}_2\text{O}$)
    \item $6\text{ electrons}$
\end{enumerate}

\item In **Complex I ($\text{NADH}$ dehydrogenase)**, electrons from $\text{NADH} + \text{H}^+$ are first transferred to which primary prosthetic group before passing through iron-sulfur ($\text{Fe}-\text{S}$) clusters to ubiquinone?
\begin{enumerate}[label=(\Alph*)]
    \item Flavin adenine dinucleotide ($\text{FAD}$)
    \item Flavin mononucleotide ($\text{FMN}$)
    \item Heme $b$
    \item Copper ($\text{Cu}_A$)
\end{enumerate}

\item Which of the following statements about **Ubiquinone** is **false**?
\begin{enumerate}[label=(\Alph*)]
    \item It is located within the inner mitochondrial membrane.
    \item It receives reducing equivalents via $\text{FADH}_2$ generated during the oxidation of succinate in the citric acid cycle.
    \item It also receives electrons from Complex I ($\text{NADH}$ dehydrogenase).
    \item The reduced ubiquinone (ubiquinol) is oxidized with the transfer of electrons directly to Cytochrome $c$, bypassing Complex III.
\end{enumerate}

\item If a mutation renders **Complex III (Cytochrome $bc_1$ complex)** completely non-functional, what will be the redox state of **Ubiquinone** and **Cytochrome $c$**, respectively, in the presence of substrates and oxygen?
\begin{enumerate}[label=(\Alph*)]
    \item Ubiquinone will be fully oxidized; Cytochrome $c$ will be fully reduced.
    \item Ubiquinone will be fully reduced (as ubiquinol, $\text{QH}_2$); Cytochrome $c$ will be fully oxidized.
    \item Both Ubiquinone and Cytochrome $c$ will be fully reduced.
    \item Both Ubiquinone and Cytochrome $c$ will be fully oxidized.
\end{enumerate}

\item In **Paul Boyer's Binding-Change Mechanism** for $\text{ATP}$ synthase ($\text{F}_1$ headpiece), the three catalytic $\beta$-subunits cycle sequentially through three conformational states ($\text{O}$, $\text{L}$, and $\text{T}$). What is the exact function of proton flow through $\text{F}_0$ in this mechanism?
\begin{enumerate}[label=(\Alph*)]
    \item Proton flow directly phosphorylates $\text{ADP}$ by forming a covalent $\text{H}^+-\text{ADP}$ bond.
    \item Proton flow drives rotation of the $c$-ring and $\gamma$-shaft, providing mechanical energy primarily to **release tightly bound, newly synthesized $\text{ATP}$** from the $\text{T}$ (tight) site by converting it into the $\text{O}$ (open) state.
    \item Proton flow cleaves $\text{ATP}$ into $\text{ADP}$ and $\text{P}_i$ to generate heat.
    \item Proton flow transports $\text{ATP}$ out of the cell into the bloodstream.
\end{enumerate}

\item Which inhibitor binds specifically to the **$\text{Q}_i$ site of Complex III (Cytochrome $bc_1$ complex)**, blocking electron transfer from cytochrome $b_H$ to ubiquinone?
\begin{enumerate}[label=(\Alph*)]
    \item Rotenone
    \item Antimycin A
    \item Malonate
    \item Cyanide
\end{enumerate}

\item Calculate the total number of $\text{ATP}$ molecules generated **exclusively by oxidative phosphorylation** (via ETS and Complex V) during the complete aerobic oxidation of **one molecule of glucose** (using NCERT's $3\text{ ATP}/\text{NADH}$ and $2\text{ ATP}/\text{FADH}_2$, with malate-aspartate shuttle):
\begin{enumerate}[label=(\Alph*)]
    \item $4\text{ ATP}$
    \item $32\text{ ATP}$
    \item $34\text{ ATP}$
    \item $38\text{ ATP}$
\end{enumerate}

\item \textbf{Statement I:} The inner mitochondrial membrane is impermeable to almost all ions and polar molecules, including $\text{H}^+$, $\text{ATP}$, $\text{ADP}$, and pyruvate, requiring specific carrier proteins. \\
\textbf{Statement II:} The outer mitochondrial membrane is freely permeable to small molecules and ions up to $\sim 5000\text{ Da}$ due to the presence of abundant **porin (VDAC)** channels.
\begin{enumerate}[label=(\Alph*)]
    \item Both Statement I and Statement II are correct.
    \item Both Statement I and Statement II are incorrect.
    \item Statement I is correct, but Statement II is incorrect.
    \item Statement I is incorrect, but Statement II is correct.
\end{enumerate}

\item How does **$\text{ATP}/\text{ADP}$ translocase** (adenine nucleotide translocase) transport $\text{ATP}^{4-}$ out of the mitochondrial matrix and $\text{ADP}^{3-}$ into the matrix across the inner membrane?
\begin{enumerate}[label=(\Alph*)]
    \item By simple passive diffusion through the lipid bilayer.
    \item By electrogenic $1:1$ antiport driven by the membrane potential ($\Delta \psi$, matrix negative), which effectively consumes the equivalent of $\sim 1\text{ H}^+$ per $\text{ATP}$ exported.
    \item By symport with $\text{Na}^+$ ions using $\text{Na}^+/\text{K}^+$-ATPase.
    \item By endocytosis of the inner membrane.
\end{enumerate}

\item In **respiratory control** (acceptor control), what is the primary factor that regulates the rate of electron transport and oxygen consumption in healthy, tightly coupled mitochondria?
\begin{enumerate}[label=(\Alph*)]
    \item The concentration of molecular oxygen (as long as $\text{O}_2$ is above trace micromolar levels).
    \item The intracellular availability of **$\text{ADP}$ and $\text{P}_i$** (when $\text{ADP}$ rises during cellular work, $\text{ATP}$ synthase speeds up, draining the $\text{H}^+$ gradient and accelerating electron transport).
    \item The concentration of glucose in the extracellular fluid.
    \item The temperature of the outer mitochondrial membrane.
\end{enumerate}

\item Which of the following cytochromes can react directly with **molecular oxygen ($\text{O}_2$)** and is inhibited by **carbon monoxide ($\text{CO}$)** in a light-reversible manner?
\begin{enumerate}[label=(\Alph*)]
    \item Cytochrome $b$
    \item Cytochrome $c_1$
    \item Cytochrome $c$
    \item Cytochrome $a_3$
\end{enumerate}

\item Suppose isolated mitochondria are supplied with **$\beta$-hydroxybutyrate** (which generates matrix $\text{NADH}$) and **succinate** (which generates $\text{FADH}_2$). If **Rotenone** is added to the medium, what will happen to the oxidation of these two substrates?
\begin{enumerate}[label=(\Alph*)]
    \item Oxidation of both $\beta$-hydroxybutyrate and succinate will be completely blocked.
    \item Oxidation of $\beta$-hydroxybutyrate (via Complex I) will be blocked, whereas oxidation of succinate (via Complex II) will continue uninterrupted.
    \item Oxidation of succinate will be blocked, whereas $\beta$-hydroxybutyrate oxidation will continue.
    \item Both substrates will be oxidized twice as fast.
\end{enumerate}

\item Why does the **cytochrome $c$** protein remain highly conserved across evolutionarily distant eukaryotes (from yeast to plants to humans)?
\begin{enumerate}[label=(\Alph*)]
    \item Because it is encoded exclusively by nuclear introns.
    \item Because it must precisely dock with both the cytochrome $bc_1$ complex (Complex III) and cytochrome $c$ oxidase (Complex IV) via conserved positively charged lysine residues around its heme crevice.
    \item Because it synthesizes $\text{ATP}$ directly in the cytosol.
    \item Because it binds oxygen in muscle tissue like myoglobin.
\end{enumerate}

\item How many total **protons ($\text{H}^+$)** must pass through the $\text{F}_0$ particle of $\text{ATP}$ synthase to synthesize the $\text{ATP}$ produced solely by oxidative phosphorylation from **one molecule of glucose** (using NCERT's $2\text{ H}^+$ per $\text{ATP}$ and $34\text{ ATP}$ from oxidative phosphorylation)?
\begin{enumerate}[label=(\Alph*)]
    \item $34\text{ H}^+$
    \item $68\text{ H}^+$
    \item $76\text{ H}^+$
    \item $100\text{ H}^+$
\end{enumerate}

\item Which of the following is **not** a component of the mitochondrial electron transport chain?
\begin{enumerate}[label=(\Alph*)]
    \item Flavin mononucleotide ($\text{FMN}$)
    \item Iron-sulfur ($\text{Fe}-\text{S}$) proteins
    \item Plastocyanin ($\text{PC}$)
    \item Ubiquinone ($\text{CoQ}_{10}$)
\end{enumerate}

\item In the mitochondrial ETS, what is the role of **Iron-Sulfur ($\text{Fe}-\text{S}$) centers**?
\begin{enumerate}[label=(\Alph*)]
    \item They carry two electrons and two protons simultaneously across the membrane.
    \item They act as one-electron carriers in which an iron atom undergoes reversible oxidation-reduction between $\text{Fe}^{3+}$ and $\text{Fe}^{2+}$.
    \item They catalyze the decarboxylation of pyruvate.
    \item They form the transmembrane proton pore of $\text{ATP}$ synthase.
\end{enumerate}

\item \textbf{Assertion (A):} In aerobic respiration, the energy of oxidation-reduction is utilized for the production of a proton gradient required for phosphorylation. \\
\textbf{Reason (R):} Unlike photophosphorylation where light energy drives electron transport, in respiration the chemical bond energy released during sequential exergonic electron transfers from $\text{NADH}/\text{FADH}_2$ to $\text{O}_2$ powers proton translocation.
\begin{enumerate}[label=(\Alph*)]
    \item Both (A) and (R) are true and (R) is the correct explanation of (A).
    \item Both (A) and (R) are true but (R) is not the correct explanation of (A).
    \item (A) is true but (R) is false.
    \item Both (A) and (R) are false.
\end{enumerate}

\item What is the **P/O ratio** (number of molecules of $\text{ADP}$ phosphorylated to $\text{ATP}$ per atom of oxygen, $\frac{1}{2}\text{O}_2$, reduced to $\text{H}_2\text{O}$) according to classical NCERT stoichiometry when $\text{NADH} + \text{H}^+$ is oxidized vs. when $\text{FADH}_2$ is oxidized?
\begin{enumerate}[label=(\Alph*)]
    \item $3$ for $\text{NADH} + \text{H}^+$ and $2$ for $\text{FADH}_2$
    \item $2$ for $\text{NADH} + \text{H}^+$ and $3$ for $\text{FADH}_2$
    \item $38$ for $\text{NADH} + \text{H}^+$ and $2$ for $\text{FADH}_2$
    \item $1$ for $\text{NADH} + \text{H}^+$ and $1$ for $\text{FADH}_2$
\end{enumerate}

\item Which of the following statements regarding **mitochondrial DNA (mtDNA)** and the biogenesis of the ETS complexes is accurate?
\begin{enumerate}[label=(\Alph*)]
    \item All subunits of Complexes I, II, III, IV, and V are encoded exclusively by nuclear DNA.
    \item Complexes I, III, IV, and V are dual-origin complexes containing subunits encoded by both nuclear DNA and mitochondrial DNA, whereas **Complex II (Succinate dehydrogenase)** is encoded **entirely by nuclear DNA**.
    \item All subunits of the ETS are encoded exclusively by mitochondrial DNA.
    \item Mitochondria lack ribosomes and cannot translate any ETS proteins.
\end{enumerate}

\end{enumerate}

\newpage

\section*{Answer Key}

\begin{multicols}{5}
\begin{enumerate}[label=\textbf{\arabic*.}]
    \item (B)
    \item (A)
    \item (A)
    \item (C)
    \item (B)
    \item (A)
    \item (B)
    \item (B)
    \item (A)
    \item (B)
    \item (D)
    \item (B)
    \item (B)
    \item (A)
    \item (A)
    \item (B)
    \item (A)
    \item (B)
    \item (B)
    \item (B)
    \item (B)
    \item (B)
    \item (A)
    \item (B)
    \item (A)
    \item (D)
    \item (C)
    \item (B)
    \item (D)
    \item (B)
    \item (B)
    \item (B)
    \item (C)
    \item (A)
    \item (B)
    \item (B)
    \item (D)
    \item (B)
    \item (B)
    \item (B)
    \item (C)
    \item (B)
    \item (A)
    \item (A)
    \item (B)
\end{enumerate}
\end{multicols}

\section*{Detailed Solutions}

\begin{enumerate}[label=\textbf{\arabic*.}]

\item \textbf{(B)} Per NCERT Class 11 Biology:
\begin{itemize}
    \item Complex I: $\text{NADH}$ dehydrogenase
    \item Complex II: Succinate dehydrogenase
    \item Complex III: Cytochrome $bc_1$ complex
    \item Complex IV: Cytochrome $c$ oxidase (containing cytochromes $a$ and $a_3$, and two copper centers)
\end{itemize}

\item \textbf{(A)} **Ubiquinone** is a hydrophobic benzoquinone derivative mobile within the lipid core of the inner mitochondrial membrane. **Cytochrome $c$** is a small protein attached to the outer surface of the inner membrane, acting as a mobile carrier between Complex III and Complex IV.

\item \textbf{(A)} $\text{NADH} + \text{H}^+$ donates electrons to Complex I, pumping protons at Complexes I, III, and IV ($3\text{ ATP}$). $\text{FADH}_2$ donates electrons at Complex II directly to ubiquinone, bypassing Complex I's proton-pumping site, so protons are pumped only at Complexes III and IV ($2\text{ ATP}$).

\item \textbf{(C)} According to NCERT, **Complex IV (Cytochrome $c$ oxidase)** contains **cytochromes $a$ and $a_3$**, and **two copper centers** ($\text{Cu}_A$ and $\text{Cu}_B$).

\item \textbf{(B)} Both processes rely on chemiosmosis, but they differ in their primary energy source: in oxidative phosphorylation, the energy of oxidation-reduction reactions generates the proton gradient, whereas in photophosphorylation, light energy drives electron transport to establish the proton gradient.

\item \textbf{(A)} Cyanide and $\text{CO}$ bind to Complex IV, blocking terminal electron transfer to $\text{O}_2$ (P-2). Rotenone blocks Complex I $\rightarrow$ CoQ transfer (Q-3). Oligomycin plugs the $\text{F}_0$ proton channel of $\text{ATP}$ synthase (R-1). $\text{2,4-DNP}$ is a lipophilic protonophore that ferries $\text{H}^+$ across the inner membrane, uncoupling oxidation from phosphorylation (S-4).

\item \textbf{(B)} During electron transport, Complexes I, III, and IV pump protons ($\text{H}^+$) from the **mitochondrial matrix** across the inner membrane into the **intermembrane space**, creating an acidic, positively charged intermembrane space relative to the alkaline, negatively charged matrix.

\item \textbf{(B)} Verbatim from NCERT Class 11 Biology (Respiration in Plants): \textit{``For each ATP produced, $2\text{H}^+$ passes through $\text{F}_0$ from the intermembrane space to the matrix down the electrochemical proton gradient.''} (Note: Modern structural biology shows $\sim 3\text{--}4\text{ H}^+$, but NEET strictly follows the NCERT figure of **$2\text{H}^+$**).

\item \textbf{(A)} Statements (i), (ii), and (iii) are correct. Statement (iv) is incorrect because $\text{F}_0$ is embedded in the inner mitochondrial membrane and $\text{F}_1$ protrudes into the **mitochondrial matrix** (not the intermembrane space).

\item \textbf{(B)} Oxygen is essential only at the final step of aerobic respiration: it acts as the **final hydrogen/electron acceptor** at Complex IV, combining with electrons and matrix protons to form metabolic water ($\frac{1}{2}\text{O}_2 + 2e^- + 2\text{H}^+ \rightarrow \text{H}_2\text{O}$).

\item \textbf{(D)} Complete aerobic oxidation of $1\text{ glucose}$ (with malate-aspartate shuttle) produces $10\text{ NADH} + \text{H}^+$ and $2\text{ FADH}_2$ — a total of **$12\text{ pairs}$ of electrons** ($24\text{ electrons}$). Since $2\text{ electrons}$ reduce $\frac{1}{2}\text{O}_2$ to $1\text{ H}_2\text{O}$, $12\text{ electron pairs}$ form **$12\text{ molecules}$ of $\text{H}_2\text{O}$** at Complex IV.

\item \textbf{(B)} Ubiquinone ($\text{Q}$) accepts two electrons and two protons ($2e^- + 2\text{H}^+$) to become ubiquinol ($\text{QH}_2$). At Complex III, it releases the $2\text{ H}^+$ into the intermembrane space while transferring electrons one at a time to cytochromes, functioning as a crucial **proton-electron coupler**.

\item \textbf{(B)} Uncouplers like $\text{2,4-DNP}$ collapse the proton gradient by carrying $\text{H}^+$ directly across the inner membrane. Because there is no proton back-pressure, electron transport and $\text{O}_2$ consumption run at maximum speed, but $\text{ATP}$ synthase cannot function, so all free energy is released as **heat**.

\item \textbf{(A)} Both statements are accurate. Brown adipose tissue contains **Thermogenin (UCP1)**, which uncouples oxidative phosphorylation to generate body heat (non-shivering thermogenesis) in newborns and hibernating animals.

\item \textbf{(A)} The thermodynamically downhill order of carriers in the ETS is:
$$\text{NADH} \rightarrow \text{FMN} \rightarrow \text{Fe-S} \rightarrow \text{Ubiquinone} \rightarrow \text{Cyt } b \rightarrow \text{Cyt } c_1 \rightarrow \text{Cyt } c \rightarrow \text{Cyt } a \rightarrow \text{Cyt } a_3 \rightarrow \text{O}_2$$

\item \textbf{(B)} Complexes I, III, and IV are transmembrane proton pumps. **Complex II (Succinate dehydrogenase)** does not span the entire membrane in a proton-pumping configuration and pumps **$0\text{ protons}$**.

\item \textbf{(A)} In liver, heart, and kidney, the **malate-aspartate shuttle** regenerates matrix $\text{NADH}$ ($3\text{ ATP}$ per cytosolic $\text{NADH}$, total $38\text{ ATP}$/glucose). In skeletal muscle and brain, the **glycerol-3-phosphate shuttle** transfers electrons to mitochondrial $\text{FAD}$ ($2\text{ ATP}$ per cytosolic $\text{NADH}$, total $36\text{ ATP}$/glucose).

\item \textbf{(B)} In skeletal muscle (glycerol-3-phosphate shuttle), the $2\text{ cytosolic NADH}$ from glycolysis yield $2 \times 2 = 4\text{ ATP}$ instead of $6\text{ ATP}$. Thus, total yield = $38 - 2 = \mathbf{36\text{ ATP}}$.

\item \textbf{(B)} Through the **Q-cycle** at Complex III, ubiquinol ($\text{QH}_2$) splits its two electrons: one reduces the Rieske $\text{Fe}-\text{S}$ center $\rightarrow$ cytochrome $c_1$ $\rightarrow$ cytochrome $c$, while the second travels through cytochrome $b_L$ and $b_H$ to partially reduce an oxidized ubiquinone to a semiquinone radical.

\item \textbf{(B)} Cyanide ($\text{CN}^-$) binds with extremely high affinity to the ferric ($\text{Fe}^{3+}$) heme $a_3$ of **Complex IV**, completely blocking electron transfer to $\text{O}_2$ and causing immediate cellular energy failure (histotoxic hypoxia).

\item \textbf{(B)} In intact mitochondria, oxidation and phosphorylation are **obligately coupled**. Blocking $\text{F}_0$ with oligomycin prevents protons from re-entering the matrix; the resulting steep proton-motive force exerts thermodynamic back-pressure on Complexes I, III, and IV, arresting electron flow and $\text{O}_2$ uptake.

\item \textbf{(B)} Adding $\text{2,4-DNP}$ after oligomycin provides an alternative route for $\text{H}^+$ ions to leak back into the matrix, relieving the thermodynamic back-pressure and restoring $\text{O}_2$ consumption, while $\text{ATP}$ synthesis remains blocked by oligomycin.

\item \textbf{(A)} All cytochromes ($b$, $c_1$, $c$, $a$, $a_3$) are heme-containing proteins in which an iron ion undergoes reversible one-electron oxidation-reduction ($\text{Fe}^{3+} + e^- \leftrightarrow \text{Fe}^{2+}$).

\item \textbf{(B)} As stated in NCERT, **Cytochrome $c$** is a small protein attached to the **outer surface of the inner membrane** and acts as a mobile carrier for transfer of electrons between Complex III and Complex IV.

\item \textbf{(A)} Pumping positively charged $\text{H}^+$ ions from the matrix to the intermembrane space simultaneously generates a chemical pH difference ($\Delta \text{pH}$, matrix alkaline) and an electrical voltage ($\Delta \psi$, matrix negative), together forming the proton-motive force.

\item \textbf{(D)} All four assumptions are explicitly listed in NCERT Class 11 Biology under \textit{The Respiratory Balance Sheet} as theoretical idealizations that do not strictly hold simultaneously in living cells.

\item \textbf{(C)} Complete reduction of one molecule of diatomic oxygen ($\text{O}_2$) requires **$4\text{ electrons}$** and **$4\text{ H}^+$** to yield $2\text{ molecules}$ of $\text{H}_2\text{O}$ ($\text{O}_2 + 4e^- + 4\text{H}^+ \rightarrow 2\text{H}_2\text{O}$). Complex IV holds the partially reduced oxygen intermediates tightly between heme $a_3$ and $\text{Cu}_B$ until all $4\text{ electrons}$ are delivered.

\item \textbf{(B)} In Complex I ($\text{NADH}$ dehydrogenase), electrons from $\text{NADH} + \text{H}^+$ are first accepted by **Flavin mononucleotide ($\text{FMN}$)**, reduced to $\text{FMNH}_2$, and then passed through a series of $\text{Fe}-\text{S}$ clusters to ubiquinone.

\item \textbf{(D)} Statement (D) is false because reduced ubiquinone (ubiquinol) transfers its electrons to **Complex III (Cytochrome $bc_1$ complex)**, which then passes them to Cytochrome $c$ — ubiquinol never transfers electrons directly to Cytochrome $c$.

\item \textbf{(B)} Blocking Complex III creates a dam in the electron flow: carriers upstream of the block (Complex I, Complex II, and **Ubiquinone**) become backed up with electrons (**fully reduced**), while carriers downstream of the block (**Cytochrome $c$**, Complex IV) donate their remaining electrons to $\text{O}_2$ and remain **fully oxidized**.

\item \textbf{(B)} According to Boyer's Binding-Change Mechanism, $\text{ADP} + \text{P}_i$ spontaneously condense into tightly bound $\text{ATP}$ at the $\text{T}$ site; the mechanical energy of proton-driven $\gamma$-subunit rotation is primarily required to **force open the $\text{T}$ site and release the synthesized $\text{ATP}$**.

\item \textbf{(B)} **Antimycin A** specifically inhibits Complex III (Cytochrome $bc_1$ complex) by binding to the $\text{Q}_i$ site and blocking electron flow from cytochrome $b_H$ to ubiquinone.

\item \textbf{(C)} Total $\text{ATP}$ from complete glucose oxidation = $38\text{ ATP}$. Of these, **$4\text{ ATP}$** come from substrate-level phosphorylation ($2$ net in glycolysis + $2\text{ GTP}/\text{ATP}$ in TCA cycle), and **$34\text{ ATP}$** come from oxidative phosphorylation ($10\text{ NADH} \times 3 + 2\text{ FADH}_2 \times 2 = 30 + 4 = 34\text{ ATP}$).

\item \textbf{(A)} Both statements are correct: the outer membrane has wide **porin** channels making it permeable to small solutes, whereas the inner membrane is a tight permeability barrier requiring specific antiporters/symporters so that the proton gradient is not dissipated.

\item \textbf{(B)} **$\text{ATP}/\text{ADP}$ translocase** exchanges matrix $\text{ATP}^{4-}$ for intermembrane $\text{ADP}^{3-}$ in a $1:1$ antiport. Because a net negative charge ($-1$) moves out of the negatively charged matrix, this transport is driven by $\Delta \psi$ and consumes energy equivalent to $\sim 1\text{ H}^+$.

\item \textbf{(B)} In coupled mitochondria, the rate of respiration is governed by the availability of **$\text{ADP}$ and $\text{P}_i$** (respiratory control). When cellular $\text{ATP}$ consumption rises, $[\text{ADP}]$ increases, stimulating $\text{ATP}$ synthase, lowering the proton gradient, and accelerating electron transport.

\item \textbf{(D)} **Cytochrome $a_3$** (together with $\text{Cu}_B$ in Complex IV) is the only cytochrome capable of binding and reducing molecular oxygen ($\text{O}_2$), and it also binds carbon monoxide ($\text{CO}$) competitively with $\text{O}_2$.

\item \textbf{(B)} **Rotenone** selectively inhibits Complex I ($\text{NADH}$ dehydrogenase), blocking oxidation of $\text{NADH}$-linked substrates like $\beta$-hydroxybutyrate. However, succinate enters at **Complex II**, bypassing Complex I, so its oxidation continues normally.

\item \textbf{(B)} Cytochrome $c$ is a soluble shuttle that must physically bind to both Complex III (to accept electrons) and Complex IV (to donate electrons); mutations in its conserved lysine-rich docking surface disrupt both interactions, explaining its extreme evolutionary conservation.

\item \textbf{(B)} Using NCERT's values: $34\text{ ATP}$ are produced via oxidative phosphorylation per glucose, and **$2\text{ H}^+$** pass through $\text{F}_0$ per $\text{ATP}$ formed. Therefore, $34 \times 2 = \mathbf{68\text{ H}^+}$ pass through $\text{F}_0$.

\item \textbf{(C)} **Plastocyanin ($\text{PC}$)** is a copper-containing mobile electron carrier in the **chloroplast thylakoid lumen** (photosynthetic electron transport), not in mitochondria.

\item \textbf{(B)} **Iron-sulfur ($\text{Fe}-\text{S}$) centers** (non-heme iron proteins) are single-electron carriers in which an iron atom cycles between the oxidized ferric ($\text{Fe}^{3+}$) and reduced ferrous ($\text{Fe}^{2+}$) states.

\item \textbf{(A)} Both statements are accurate and aligned with NCERT: sequential exergonic redox reactions along the ETS release free energy used by Complexes I, III, and IV to pump protons across the inner membrane, creating the electrochemical proton gradient for $\text{ATP}$ synthesis.

\item \textbf{(A)} By NCERT stoichiometry, oxidizing $1\text{ NADH} + \text{H}^+$ reduces $\frac{1}{2}\text{O}_2$ ($1\text{ oxygen atom}$) and yields **$3\text{ ATP}$** ($\text{P/O ratio} = 3$), whereas oxidizing $1\text{ FADH}_2$ reduces $\frac{1}{2}\text{O}_2$ and yields **$2\text{ ATP}$** ($\text{P/O ratio} = 2$).

\item \textbf{(B)} Mitochondria are semi-autonomous organelles: Complexes I, III, IV, and V contain core catalytic subunits encoded by **mtDNA** alongside nuclear-encoded subunits, whereas **Complex II (Succinate dehydrogenase)** is the only ETS complex encoded **100% by nuclear DNA**.

\end{enumerate}

\end{document}
</USER_REQUEST>
<ADDITIONAL_METADATA>
The current local time is: 2026-11-16T12:03:56+05:30.

The user's current state is as follows:
Active Document: /Users/laxmikumari/Desktop/web/testseries/scripts/add_botany_zoology_to_part_11.mjs (LANGUAGE_UNSPECIFIED)
Cursor is on line: 1
Other open documents:
- /Users/laxmikumari/Desktop/web/testseries/.env.local (LANGUAGE_UNSPECIFIED)
- /Users/laxmikumari/Downloads/neet_five_kingdom_classification_test.tex (LANGUAGE_UNSPECIFIED)
- /Users/laxmikumari/Downloads/GOC_Extracted_Questions.tex (LANGUAGE_UNSPECIFIED)
- /Users/laxmikumari/Downloads/neet_biology_question_bank_phases_of_growth_set_2.tex (LANGUAGE_UNSPECIFIED)
- /Users/laxmikumari/Desktop/web/testseries/src/app/api/admin/notifications/route.js (LANGUAGE_UNSPECIFIED)
</ADDITIONAL_METADATA>